% Primitive: inclusion-exclusion — authored from
% probability/inclusion_exclusion_manim.py, its README concepts table (seven
% scenes) and plan 016's verified anchors (verifier report A–K plus the
% addendum L–R). Delivers the counting queue's "two and three overlapping
% sets" from the probability side and climbs to four sets and n. Problem
% answers are computed by answers/inclusion_exclusion.py (the solve-gate
% anchor, 8/8 on a fresh-context solve); every exact value there is
% Fraction arithmetic, floats only where a decimal display is the claim.
\chapter{Inclusion–exclusion}

\begin{narrative}
The product rule answers ``$A$ \emph{and} $B$''. This chapter is its
sibling: ``$A$ \emph{or} $B$'', and then ``at least one of three, four,
$n$''. The whole subject is one sentence — \emph{whatever you counted
twice, subtract once} — and the chapter's work is to show that the
sentence, applied honestly, forces the alternating signs, survives the
picture breaking at four sets, and turns into a bound when you stop
early.

\section{Two sets, one overlap}

Roll a die; let $A$ be even and $B$ be at most four — the pair the
independence chapter already used. Tint $A$'s three cells, lay $B$'s
four over them, and two cells carry both colours. Count the union and
you get five cells, not seven: $|A \cup B| = 3 + 4 - 2$. The overlap
was counted twice, so it is subtracted once. In fractions the mistake
announces itself before any picture explains it: $P(A) + P(B) =
\anchor{016.L.alarm}$, a number no probability can be. The same
correction as area on the unit square — a band of width $1/2$, a band
of height $2/3$, the doubly covered rectangle removed once — gives
\[
  P(A \cup B) \;=\; P(A) + P(B) - P(A \cap B) \;=\; \anchor{016.L.dieunion} ,
\]
Ross's Proposition~4.3, proved from additivity by cutting $A \cup B$
into disjoint pieces~\cite{probability-sheldon-ross-a-first-course-in-probability-10th-ed-pearson-2}.
Three readings of the overlap close the scene. Disjoint events are the
case where the sum is \emph{exact} ($\{1,2\}$ and $\{5,6\}$: $2/6 +
2/6$); independent events are the case where the overlap is a
\emph{product} ($1/3 = 1/2 \cdot 2/3$, so the complement route $1 -
\tfrac12\cdot\tfrac13$ agrees) — disjoint and independent are opposite
ends of the overlap, never the same thing; and ``exactly one'' is the
same sum with the overlap subtracted \emph{twice} ($1/2$, not $5/6$).
On the two-dice grid one cell carries both sixes:
$P(\text{a six somewhere}) = \anchor{016.B.twodice}$, and $12/36$ is
the overcount, not an answer.

\section{Three sets, one ledger}

Three events on the same grid — first die six, second die six, sum at
least ten. Count the union directly: twelve cells,
$\anchor{016.M.threesets}$ (here $1/3$ \emph{is} the answer). Now
watch the formula earn it. Stamp every cell with a running count as
the terms arrive: $+1$ for each event it lies in, then $-1$ for each
pair, then $+1$ for the triple. The cell $(6,6)$ sits in all three. It
reads $3$ after the additions, $0$ after the three subtractions — it
has vanished from a union it belongs to — and $1$ after the last term.
Every cell in the union ends at $1$, every cell outside at $0$, and
the sum of the stamps is the count: $6 + 6 + 6 - 1 - 3 - 3 + 1 = 12$.
That is why the triple is \emph{added back}, not subtracted: the last
term is the only one that knows about the centre. One region of the
picture — in $A$ and $B$ but not $C$ — is empty; the formula does not
mind, which is the first hint that the picture is not doing the work.

Bernstein's coins from the independence chapter (first heads, second
heads, exactly one head) make the point twice
over~\cite{probability-alexander-bogomolny-mutually-jointly-independent-events}.
Their triple intersection is empty, so $P(A \cup B \cup C) =
\tfrac32 - \tfrac34 + 0 = \anchor{016.N.bernstein}$. And the shortcut
``multiply the complements'' returns $1 - (1/2)^3 = 7/8$: every pair
term was a correct product, the triple was not, and the error is
exactly the product the true triple term is not. Pairwise independence
licenses the pairs; only \emph{mutual} independence licenses the
shortcut~\cite{probability-wikipedia-pairwise-independence}.

\section{Four sets, no picture}

Draw four overlapping circles and number the regions. You reach
fourteen — thirteen inside plus the outside — against the sixteen a
four-set Venn diagram needs; the two missing are the ``opposite pairs
only'' regions. Venn knew: ``four circles cannot be so drawn as to
intersect one another in the way required'', and supplied four
ellipses that do all sixteen~\cite{probability-john-venn-symbolic-logic-macmillan-1881};
the count $\anchor{016.F.circles}$ is Euler's formula with twelve
crossing points and twenty-four
arcs~\cite{probability-satyadev-nandakumar-venn-diagrams-and-circles-iit-kanpur}.
The honest conclusion is not that ellipses save the picture but that
the picture was never area-true — the grid was — and the ledger needs
no picture at all.

Two four-set examples, one coefficient row. Four hats handed back at
random: the nine derangements of $1234$ are the permutations with no
hat home, and $P(\text{at least one match}) = 4\cdot\tfrac14 -
6\cdot\tfrac1{12} + 4\cdot\tfrac1{24} - \tfrac1{24} =
\anchor{016.E.fourhats}$, with no product anywhere. De Méré's four
rolls of a die: $P(\text{at least one six}) = 4\cdot\tfrac16 -
6\cdot\tfrac1{36} + 4\cdot\tfrac1{216} - \tfrac1{1296} =
\anchor{016.O.fourrolls}$, every term a
product~\cite{probability-eric-w-weisstein-de-m-r-s-problem-mathworld}.
The coefficients $4, 6, 4, 1$ are the same both times — the sorted
square's column counts from the random-variables chapter — and any
outcome lying in all four events is counted $\anchor{016.G.fourledger}$.
When every $k$-fold intersection has the same probability $p_k$ the
whole sum collapses to $\sum_k (-1)^{k+1}\binom{n}{k} p_k$; the
complement $1 - (5/6)^4$ says the same thing in one line for the
rolls, and for the hats no such line exists.

\section{Every point counted once}

For $n$ events the formula has one term per non-empty subset $S$,
its sign alternating with $|S|$:
\[
  P\Big(\bigcup_{i=1}^{n} A_i\Big)
  \;=\; \sum_{k=1}^{n} (-1)^{k+1} \sum_{|S| = k}
        P\Big(\bigcap_{i \in S} A_i\Big) ,
\]
$2^n - 1$ terms in all~\cite{probability-william-feller-an-introduction-to-probability-theory-vol-1-1}.
The reason it is right is the ledger, said once for all $n$: an
outcome lying in exactly $k$ of the events is counted $\binom{k}{1} -
\binom{k}{2} + \binom{k}{3} - \cdots$ times, and that alternating sum
is $1$. The textbooks recall $(1-1)^k = 0$ for this — the binomial
theorem, which this road has not built — so the chapter proves it by
pairing instead~\cite{probability-dennis-white-math-4707-inclusion-exclusion-and-derangements,probability-benjamin-and-quinn-an-alternate-approach-to-alternating-sums}.
List the non-empty subsets of $\{1, \dots, k\}$ and pair each $S$
with $S \,\triangle\, \{1\}$ (toggle whether $1$ is a member). Every
pair has sizes of opposite parity, so its two terms cancel; the only
subset without a partner is $\{1\}$ itself, and its term is $+1$. So
$3 - 3 + 1 = 1$, $4 - 6 + 4 - 1 = 1$, and so on for every $k$. The
running partial sums ($4, -2, 2, 1$) are counts, not probabilities —
a cell can be counted $-2$ times mid-ledger — and their sign pattern
is the next-but-one section's tool.

\section{The matching limit}

Return $n$ hats at random. With $\binom{n}{k}$ equal $k$-fold terms
each worth $(n-k)!/n!$, the coefficient cancels the count and
\[
  P(\text{at least one match}) \;=\;
  1 - \frac{1}{2!} + \frac{1}{3!} - \cdots \pm \frac{1}{n!} ,
\]
the values $\tfrac12, \tfrac23, \tfrac58, \tfrac{19}{30}, \dots$
oscillating — above the limit at $n = 3$, below at $n = 4$ — with
each error smaller than the next term. The limit is
$1 - 1/e \approx \anchor{016.J.limit}$: more hats never make a match
certain~\cite{probability-wikipedia-derangement}. The road already
met $1/e$ once, as the binomial's zero-success limit $(1 - 1/n)^n$;
the two roads differ at every $n \ge 2$ ($0.3164$ against $0.375$ at
$n = 4$) and one of them has no independence in it. The identity
$\sum (-1)^k / k! = e^{-1}$ is a series fact this road has not built;
$e$ itself is the calculus chapter's compound-interest ceiling. The
problem is Montmort's \emph{Treize}, posed in 1708 and solved in
1713~\cite{probability-o-connor-and-robertson-montmort-s-probl-me-du-treize-mactuto};
de Moivre's 1718 \emph{Doctrine of Chances} states the general
case~\cite{probability-abraham-de-moivre-the-doctrine-of-chances-london-1718}.

\section{Brackets and bounds}

Stop the sum early and you hold a bound: one term is an upper bound
(Boole's inequality), two terms a lower bound, three an upper bound,
and so on — Ross's Remark~3~\cite{probability-sheldon-ross-a-first-course-in-probability-10th-ed-pearson-2}.
On the four hats the partial sums run $\anchor{016.I.brackets}$; on
the four rolls $\tfrac23, \tfrac12, \tfrac{14}{27}, \tfrac{671}{1296}$
— over, under, over, and the last term lands exactly. The reason is
again the ledger: a truncated sum over- or under-counts every outcome
with the sign of its last kept term, since $\sum_{j=1}^{m} (-1)^{j+1}
\binom{k}{j} = 1 - (-1)^m \binom{k-1}{m}$ (for $k = 4$: $4, -2, 2, 1$)~\cite{probability-benjamin-and-quinn-an-alternate-approach-to-alternating-sums}.
The first rung is the whole tool when events are rare: four events
at $0.01$ have union at most $0.04$, against $\anchor{016.R.rare}$
exactly under independence — and the bound needed no independence at
all. On four hats the same rung says ``at most $1$'': honest and
useless, because the union bound bites only when $\sum P(A_i) < 1$.
Bonferroni named the ladder in 1936~\cite{probability-wikipedia-boole-s-inequality}.

\section{When to use it}

Disjoint events: add, the sum is exact. Independent events, ``at
least one'': the complement $1 - \prod (1 - p_i)$. Dependent but
symmetric — hats, matches: inclusion–exclusion with
$\binom{n}{k} p_k$. Many rare events, a guarantee: the union bound,
first term only. ``Exactly $m$ of $n$'': the sieve, named here and
not built. The same rule read as a count is a sieve: among $1$ to
$30$, the integers divisible by $2$, $3$ or $5$ number $15 + 10 + 6 -
5 - 3 - 2 + 1 = \anchor{016.C.sieve}$, leaving eight survivors —
Euler's $\varphi(30)$. Blitzstein and Hwang's verdict closes the
series: try the other tools first; inclusion–exclusion is the last
resort~\cite{probability-blitzstein-hwang-introduction-to-probability-chs-1-2-excerpt}.
\end{narrative}

\section*{Practice problems}

\begin{problem}
A fair die is rolled once. Let $A$ be ``at most $3$'' and $B$ be ``a
multiple of $3$''. Find $P(A \cup B)$ and the probability that exactly
one of $A$, $B$ occurs.
\begin{hint}
Which face is in both? Count it once for the union, and not at all for
``exactly one''.
\end{hint}
\begin{solution}
$A = \{1,2,3\}$, $B = \{3,6\}$, $A \cap B = \{3\}$. The union is
$\tfrac36 + \tfrac26 - \tfrac16 = \tfrac46 = \tfrac23$. Exactly one
subtracts the overlap twice: $\tfrac36 + \tfrac26 - 2\cdot\tfrac16 =
\tfrac12$ — the faces $\{1, 2, 6\}$.
\end{solution}
\end{problem}

\begin{problem}
Two fair dice. $A$: the first die shows $5$ or $6$; $B$: the second
die shows $5$ or $6$; $C$: the faces sum to at least $11$. Give the
seven counts $|A|, |B|, |C|, |A \cap B|, |A \cap C|, |B \cap C|,
|A \cap B \cap C|$ and $P(A \cup B \cup C)$. Then run the ledger on
the outcome $(6,6)$: how many times is it counted after the three
``$+$'' terms, after the three ``$-$'' terms, and at the end?
\begin{hint}
$C$ has only three outcomes, and all of them lie in both $A$ and $B$.
\end{hint}
\begin{solution}
$|A| = |B| = 12$, $|C| = 3$ (the cells $(5,6)$, $(6,5)$, $(6,6)$);
$|A \cap B| = 4$; $A \cap C = B \cap C = A \cap B \cap C = C$, so those
three counts are $3$. The union is $12 + 12 + 3 - 4 - 3 - 3 + 3 = 20$
cells, $P = \tfrac{20}{36} = \tfrac59$. The outcome $(6,6)$ lies in all
three events: it reads $3$ after the additions, $0$ after the three
subtractions, and $1$ after the triple is added back.
\end{solution}
\end{problem}

\begin{problem}
Five hats are returned at random to their five owners. What is the
probability that at least one owner receives their own hat? How many
of the $120$ permutations have no fixed point?
\begin{hint}
Five equal terms at $k$: $\binom{5}{k}\,(5-k)!/5! = 1/k!$.
\end{hint}
\begin{solution}
$1 - \tfrac1{2!} + \tfrac1{3!} - \tfrac1{4!} + \tfrac1{5!} =
\tfrac{76}{120} = \tfrac{19}{30}$. The derangements number $120 - 76
= 44$.
\end{solution}
\end{problem}

\begin{problem}
A fair die is rolled three times. Find the probability of at least
one six by inclusion–exclusion (state the three terms) and by the
complement. Then give the partial sums after one and after two terms,
and order them against the exact value.
\begin{hint}
The events ``roll $i$ is a six'' are independent, so every
intersection is a product — and the first partial sum is already
$1/2$.
\end{hint}
\begin{solution}
$3\cdot\tfrac16 - 3\cdot\tfrac1{36} + \tfrac1{216} = \tfrac{108 - 18 +
1}{216} = \tfrac{91}{216}$; the complement $1 - (5/6)^3 = 1 -
\tfrac{125}{216}$ agrees. The partial sums are $\tfrac12 =
\tfrac{108}{216}$ and $\tfrac{90}{216} = \tfrac5{12}$, and
$\tfrac5{12} \le \tfrac{91}{216} \le \tfrac12$: over, then under, then
exact.
\end{solution}
\end{problem}

\begin{problem}
Let $\Omega = \{1, 2, 3, 4\}$ with equal probabilities, $A = \{1,2\}$,
$B = \{1,3\}$, $C = \{1,4\}$. Check each pair for independence; give
$P(A \cap B \cap C)$; compute $P(A \cup B \cup C)$ by
inclusion–exclusion; compute $1 - P(A^c)P(B^c)P(C^c)$; and give the
signed difference between that shortcut and the true union.
\begin{hint}
Every intersection — pairs and triple alike — is the single outcome
$\{1\}$.
\end{hint}
\begin{solution}
Each event has probability $\tfrac12$ and each pair meets in $\{1\}$
with probability $\tfrac14 = \tfrac12\cdot\tfrac12$: pairwise
independent. The triple is also $\{1\}$, $\tfrac14 \ne \tfrac18$: not
mutually independent. The union is $\tfrac32 - \tfrac34 + \tfrac14 =
1$ — every outcome is covered. The shortcut gives $1 - \tfrac18 =
\tfrac78$, low by $\tfrac18$: the same error as the chapter's coins,
with the opposite sign, because here the true triple term is
$\tfrac14$ and the product $\tfrac18$ is too small.
\end{solution}
\end{problem}

\begin{problem}
Among the integers $1$ to $60$, how many are divisible by $2$, $3$ or
$5$? How many by none of the three?
\begin{hint}
The same rule as a count: $\lfloor 60/2 \rfloor + \lfloor 60/3 \rfloor
+ \lfloor 60/5 \rfloor - \cdots$.
\end{hint}
\begin{solution}
$30 + 20 + 12 - 10 - 6 - 4 + 2 = 44$ are divisible by at least one;
$60 - 44 = 16$ by none — which is $60 \cdot \tfrac12 \cdot \tfrac23
\cdot \tfrac45$, the sieve in product form.
\end{solution}
\end{problem}

\begin{problem}
Take the non-empty subsets of $\{1, 2, 3, 4, 5\}$ and the map $S
\mapsto S \,\triangle\, \{1\}$. How many subsets are paired with a
different non-empty subset, how many pairs is that, and which subset
is left alone? Then evaluate $\binom51 - \binom52 + \binom53 -
\binom54 + \binom55$, listing the running partial sums.
\begin{hint}
The map toggles $1$; only one non-empty subset toggles to the empty
set.
\end{hint}
\begin{solution}
Thirty subsets form fifteen pairs; $\{1\}$ is alone, since its partner
would be $\varnothing$. Each pair holds one odd-sized and one
even-sized set, so their contributions cancel and the sum is the
survivor's $+1$: $5 - 10 + 10 - 5 + 1 = 1$, with partial sums $5, -5,
5, 0, 1$.
\end{solution}
\end{problem}

\begin{problem}[stretch]
Five events each have probability $0.02$. What does Boole's
inequality say about the probability of their union? If the events
are independent, what is that probability exactly, to six decimal
places?
\begin{hint}
The bound needs no independence; the exact value does.
\end{hint}
\begin{solution}
The union bound gives $P(\bigcup A_i) \le 5 \times 0.02 = 0.10$.
Under independence $P = 1 - 0.98^5 = 0.096079$ — the first rung is
about four per cent above the truth, and it would still hold if the
events were dependent in any way at all.
\end{solution}
\end{problem}
